% Generated by make_conventional_newest.py from the clean32 CSV files.
\newcommand{\ConvChartLabels}{%
0/{NoPF}/muted,1/{Next-Line}/muted,2/{Stride}/muted,3/{Streamer}/muted,4/{AMPM}/muted,5/{SPP}/muted,6/{Newest}/teal%
}
\newcommand{\ConvRegionHeaders}{%
0/{410M},1/{430M},2/{450M}%
}
\newcommand{\ConvIpcBars}{%
0/0/0.85145738/0.85146/muted!62,0/1/0.85545584/0.85546/muted!62,0/2/0.85426484/0.85426/muted!62,0/3/0.85820746/0.85821/muted!62,0/4/0.85749657/0.85750/muted!62,0/5/0.85842774/0.85843/muted!62,0/6/0.86322172/0.86322/teal,1/0/0.19544978/0.19545/muted!62,1/1/0.19463910/0.19464/muted!62,1/2/0.19622916/0.19623/muted!62,1/3/0.20677887/0.20678/muted!62,1/4/0.20426302/0.20426/muted!62,1/5/0.20755452/0.20755/muted!62,1/6/0.21574691/0.21575/teal,2/0/0.63862920/0.63863/muted!62,2/1/0.63281737/0.63282/muted!62,2/2/0.63857659/0.63858/muted!62,2/3/0.65189059/0.65189/muted!62,2/4/0.65110621/0.65111/muted!62,2/5/0.64972002/0.64972/muted!62,2/6/0.66369591/0.66370/teal%
}
\newcommand{\ConvPhaseTable}{%
\addlinespace[2pt]\multicolumn{12}{@{}l}{\textbf{410M} \quad NoPF misses 13,550} \\%
NoPF & 0.85146 & 1.174 & — & — & 0.00 & 0.00 & 60.93 & 39.07 & 3.387 & 0.869 & 0 \\%
Next-Line & 0.85546 & 1.169 & 42.09 & 41.01 & 43.98 & 18.01 & 67.21 & 32.79 & 2.776 & 1.394 & 14,177 \\%
Stride & 0.85426 & 1.171 & 87.32 & 97.10 & 1.78 & 1.73 & 60.86 & 39.14 & 3.329 & 0.871 & 276 \\%
Streamer & 0.85821 & 1.165 & 27.64 & 57.67 & 24.04 & 13.90 & 65.80 & 34.20 & 2.915 & 1.419 & 11,864 \\%
AMPM & 0.85750 & 1.166 & 25.01 & 87.49 & 10.56 & 9.24 & 63.77 & 36.23 & 3.074 & 1.144 & 5,721 \\%
SPP & 0.85843 & 1.165 & 44.75 & 49.26 & 24.34 & 11.99 & 65.01 & 34.99 & 2.980 & 1.131 & 7,394 \\%
\textbf{Newest} & 0.86322 & 1.158 & 76.19 & 57.75 & 28.74 & 16.64 & 66.89 & 33.11 & 2.823 & 0.947 & 5,129 \\%
\addlinespace[2pt]\multicolumn{12}{@{}l}{\textbf{430M} \quad NoPF misses 235,615} \\%
NoPF & 0.19545 & 5.116 & — & — & 0.00 & 0.00 & 0.00 & 100.00 & 58.904 & 15.237 & 0 \\%
Next-Line & 0.19464 & 5.138 & 27.53 & 41.46 & 22.18 & 9.19 & 9.20 & 90.80 & 53.483 & 24.043 & 190,005 \\%
Stride & 0.19623 & 5.096 & 92.48 & 98.55 & 0.58 & 0.58 & 0.58 & 99.42 & 58.564 & 15.225 & 1,489 \\%
Streamer & 0.20678 & 4.836 & 46.54 & 76.36 & 15.07 & 11.50 & 11.52 & 88.48 & 52.119 & 17.775 & 76,360 \\%
AMPM & 0.20426 & 4.896 & 52.58 & 96.22 & 8.31 & 8.00 & 8.00 & 92.00 & 54.188 & 16.268 & 37,272 \\%
SPP & 0.20755 & 4.818 & 45.77 & 66.72 & 16.81 & 11.21 & 11.22 & 88.78 & 52.295 & 18.158 & 86,575 \\%
\textbf{Newest} & 0.21575 & 4.635 & 42.28 & 78.28 & 24.00 & 18.79 & 18.80 & 81.20 & 47.832 & 20.180 & 133,815 \\%
\addlinespace[2pt]\multicolumn{12}{@{}l}{\textbf{450M} \quad NoPF misses 45,588} \\%
NoPF & 0.63863 & 1.566 & — & — & 0.00 & 0.00 & 46.24 & 53.76 & 11.397 & 2.923 & 0 \\%
Next-Line & 0.63282 & 1.580 & 49.43 & 35.29 & 42.82 & 15.39 & 51.74 & 48.26 & 10.041 & 4.400 & 41,594 \\%
Stride & 0.63858 & 1.566 & 84.47 & 98.61 & 0.79 & 0.78 & 45.82 & 54.18 & 11.351 & 2.938 & 425 \\%
Streamer & 0.65189 & 1.534 & 35.85 & 55.69 & 22.55 & 12.60 & 51.91 & 48.09 & 10.043 & 4.167 & 29,366 \\%
AMPM & 0.65111 & 1.536 & 32.49 & 93.25 & 11.26 & 10.52 & 50.62 & 49.38 & 10.308 & 3.654 & 16,043 \\%
SPP & 0.64972 & 1.539 & 54.67 & 46.26 & 23.50 & 10.95 & 51.43 & 48.57 & 10.041 & 3.481 & 19,912 \\%
\textbf{Newest} & 0.66370 & 1.507 & 78.20 & 57.19 & 24.57 & 14.09 & 54.36 & 45.64 & 9.571 & 3.069 & 14,418 \\%
}
\newcommand{\ConvPooledHeadA}{%
0/{IPC},1/{Time (ms)},2/{Accuracy \%},3/{Timeliness \%},4/{Coverage \%}%
}
\newcommand{\ConvPooledHeadB}{%
0/{Timely cov. \%},1/{LLC hit \%},2/{LLC miss \%},3/{MPKI},4/{DRAM (MB)}%
}
\newcommand{\ConvPooledBarsA}{%
0/0/0.92917030/{0.38184}/muted!62,0/1/0.92561216/{0.38038}/muted!62,0/2/0.93202353/{0.38301}/muted!62,0/3/0.96880314/{0.39813}/muted!62,0/4/0.96084905/{0.39486}/muted!62,0/5/0.97050932/{0.39883}/muted!62,0/6/1.00000000/{0.41095}/teal,1/0/0.99617072/{7.857}/muted!62,1/1/1.00000000/{7.887}/muted!62,1/2/0.99312111/{7.833}/muted!62,1/3/0.95541821/{7.535}/muted!62,1/4/0.96332735/{7.598}/muted!62,1/5/0.95373856/{7.522}/muted!62,1/6/0.92561224/{7.300}/teal,2/0/0.00000000/{—}/muted!62,2/1/0.32075142/{32.08}/muted!62,2/2/0.90273973/{90.27}/muted!62,2/3/0.41964453/{41.96}/muted!62,2/4/0.44452537/{44.45}/muted!62,2/5/0.47258103/{47.26}/muted!62,2/6/0.46790600/{46.79}/teal,3/0/0.00000000/{—}/muted!62,3/1/0.39813276/{39.81}/muted!62,3/2/0.98381386/{98.38}/muted!62,3/3/0.70704819/{70.70}/muted!62,3/4/0.95152993/{95.15}/muted!62,3/5/0.61511019/{61.51}/muted!62,3/6/0.73847183/{73.85}/teal,4/0/0.00000000/{0.00}/muted!62,4/1/0.26373608/{26.37}/muted!62,4/2/0.00670731/{0.67}/muted!62,4/3/0.16636981/{16.64}/muted!62,4/4/0.08870817/{8.87}/muted!62,4/5/0.18189467/{18.19}/muted!62,4/6/0.24308489/{24.31}/teal%
}
\newcommand{\ConvPooledBarsB}{%
0/0/0.00000000/{0.00}/muted!62,0/1/0.10555957/{10.56}/muted!62,0/2/0.00659875/{0.66}/muted!62,0/3/0.11784443/{11.78}/muted!62,0/4/0.08446394/{8.45}/muted!62,0/5/0.11209046/{11.21}/muted!62,0/6/0.17967247/{17.97}/teal,1/0/0.16992264/{16.99}/muted!62,1/1/0.24807054/{24.81}/muted!62,1/2/0.17105550/{17.11}/muted!62,1/3/0.26309562/{26.31}/muted!62,1/4/0.23444723/{23.44}/muted!62,1/5/0.25858091/{25.86}/muted!62,1/6/0.31873750/{31.87}/teal,2/0/0.83007736/{83.01}/muted!62,2/1/0.75192946/{75.19}/muted!62,2/2/0.82894450/{82.89}/muted!62,2/3/0.73690438/{73.69}/muted!62,2/4/0.76555277/{76.56}/muted!62,2/5/0.74141909/{74.14}/muted!62,2/6/0.68126250/{68.13}/teal,3/0/1.00000000/{24.563}/muted!62,3/1/0.89972968/{22.100}/muted!62,3/2/0.99397801/{24.415}/muted!62,3/3/0.88314630/{21.692}/muted!62,3/4/0.91698141/{22.524}/muted!62,3/5/0.88637612/{21.772}/muted!62,3/6/0.81731484/{20.075}/teal,4/0/0.63773503/{19.029}/muted!62,4/1/1.00000000/{29.838}/muted!62,4/2/0.63793451/{19.034}/muted!62,4/3/0.78290656/{23.360}/muted!62,4/4/0.70600411/{21.066}/muted!62,4/5/0.76312595/{22.770}/muted!62,4/6/0.81090443/{24.196}/teal%
}
\newcommand{\ConvPooledTable}{%
NoPF & 0.38184 & 7.857 & — & — & 0.00 & 0.00 & 16.99 & 83.01 & 24.563 & 19.029 & 0 \\%
Next-Line & 0.38038 & 7.887 & 32.08 & 39.81 & 26.37 & 10.56 & 24.81 & 75.19 & 22.100 & 29.838 & 245,776 \\%
Stride & 0.38301 & 7.833 & 90.27 & 98.38 & 0.67 & 0.66 & 17.11 & 82.89 & 24.415 & 19.034 & 2,190 \\%
Streamer & 0.39813 & 7.535 & 41.96 & 70.70 & 16.64 & 11.78 & 26.31 & 73.69 & 21.692 & 23.360 & 117,590 \\%
AMPM & 0.39486 & 7.598 & 44.45 & 95.15 & 8.87 & 8.45 & 23.44 & 76.56 & 22.524 & 21.066 & 59,036 \\%
SPP & 0.39883 & 7.522 & 47.26 & 61.51 & 18.19 & 11.21 & 25.86 & 74.14 & 21.772 & 22.770 & 113,881 \\%
\textbf{Newest} & 0.41095 & 7.300 & 46.79 & 73.85 & 24.31 & 17.97 & 31.87 & 68.13 & 20.075 & 24.196 & 153,362 \\%
}


\newcommand{\conventionaltablehead}{%
Policy & IPC & \shortstack{Time\\(ms)} & \shortstack{Accuracy\\\%}
& \shortstack{Timeliness\\\%} & \shortstack{Coverage\\\%}
& \shortstack{Timely cov.\\\%} & \shortstack{LLC hit\\\%}
& \shortstack{LLC miss\\\%} & MPKI & \shortstack{DRAM\\(MB)} & Accepted \\}

\newcommand{\conventionalpooledplot}[2]{%
\begin{tikzpicture}[x=1cm,y=.36cm,font=\footnotesize]
  \foreach \col/\heading in #1 {
    \node[font=\footnotesize\bfseries] at ({2.93*\col+1.13},.78) {\heading};
    \draw[muted!24] ({2.93*\col},.38)--({2.93*\col},-6.35);
  }
  \foreach \row/\label/\tone in \ConvChartLabels {
    \node[anchor=east,text=\tone] at (-.11,-\row) {\label};
  }
  \foreach \col/\row/\value/\raw/\tone in #2 {
    \fill[\tone] ({2.93*\col},{-\row-.25}) rectangle
        ({2.93*\col+1.43*\value},{-\row+.25});
    \node[anchor=west,font=\footnotesize] at ({2.93*\col+1.50},-\row) {\raw};
  }
\end{tikzpicture}\par}

\newpage
\sectionline{CONVENTIONAL BASELINES}
\titleline{Newest 在每个测试段跑得怎样？}
三段都是同一份 SPEC 605 MCF trace 上的不同区间。NoPF 关闭预取，给 coverage 提供分母。Next-Line 试同页下一条线，Stride 跟踪重复地址间距，Streamer 跟踪地址流，AMPM 识别页内访问图案，SPP 跟踪地址间距的序列。各策略在同一 cache 配置和每次触发至多接受 1 个请求的规则下重放。没有额外的整段发包上限。

\subhead{IPC by test region}
\begin{tikzpicture}[x=1cm,y=.43cm,font=\footnotesize]
  \foreach \col/\heading in \ConvRegionHeaders {
    \node[font=\footnotesize\bfseries] at ({5.10*\col+1.90},.91) {\heading};
    \foreach \tick in {0,0.5,1} {
      \draw[muted!25] ({5.10*\col+3.42*\tick},.42)--({5.10*\col+3.42*\tick},-6.45);
      \node[anchor=north,font=\scriptsize,text=muted] at ({5.10*\col+3.42*\tick},-6.55) {\tick};
    }
  }
  \foreach \row/\label/\tone in \ConvChartLabels {
    \node[anchor=east,text=\tone] at (-.11,-\row) {\label};
  }
  \foreach \col/\row/\value/\raw/\tone in \ConvIpcBars {
    \fill[\tone] ({5.10*\col},{-\row-.27}) rectangle
        ({5.10*\col+3.42*\value},{-\row+.27});
    \node[anchor=west,font=\footnotesize] at ({5.10*\col+3.48},-\row) {\raw};
  }
\end{tikzpicture}\par

\subhead{Complete per-region results}
{\footnotesize\setlength{\tabcolsep}{3.2pt}\renewcommand{\arraystretch}{1.08}%
\begin{tabular}{@{}lrrrrrrrrrrr@{}}
\toprule
\conventionaltablehead
\midrule
\ConvPhaseTable
\bottomrule
\end{tabular}\par}
\smallnote{这里逐段对比，保留 430M 这类运行较慢的区间。Time (ms) 是该段最后一个 core 完成所需的时间。Accepted 只数进入预取队列的请求，DRAM (MB) 数读写总字节，1 MB = 1,000,000 B。}

\newpage
\sectionline{CONVENTIONAL BASELINES}
\titleline{三段合起来，速度和流量各花了多少？}
把三段视作连续做完的三件工作。Time 加总各段的完成时间，IPC 用总退休指令数除以总完成周期。Accuracy、Coverage、LLC hit 等比例都先加总各自的分子和分母再相除。这样各段按实际事件数计入结果。

\subhead{Pooled results}
\smallnote{每格的 bar 从 0 开始。百分比刻度到 100，其余各格到本格最大值。数值标在 bar 旁边。}

\conventionalpooledplot{\ConvPooledHeadA}{\ConvPooledBarsA}
\vspace{4pt}
\conventionalpooledplot{\ConvPooledHeadB}{\ConvPooledBarsB}

{\footnotesize\setlength{\tabcolsep}{3.2pt}\renewcommand{\arraystretch}{1.12}%
\begin{tabular}{@{}lrrrrrrrrrrr@{}}
\toprule
\conventionaltablehead
\midrule
\ConvPooledTable
\bottomrule
\end{tabular}\par}

\subhead{怎么读这组结果？}
Newest 在三段都比该段最快的传统预取器快，传统最快分别是 SPP、SPP、Streamer。合计完成时间是 \textbf{7.300 ms}，逐段用最快传统策略合计是 \textbf{7.517 ms}，Newest 快 \textbf{2.88\%}，DRAM 流量多 \textbf{3.15\%}。Next-Line 的 Coverage 最高（26.37\%），Newest 的 Timely coverage 最高（17.97\%），说明被覆盖的 miss 数和提前到达的 miss 数值得分开看。
